\documentclass[11pt]{amsart}
\usepackage[T1]{fontenc}
\usepackage{lmodern}
\usepackage{amsmath,amssymb,amsthm,mathtools}
\usepackage{booktabs}
\usepackage{xcolor}
\usepackage[colorlinks=true,linkcolor=blue!55!black,citecolor=blue!55!black,urlcolor=blue!55!black]{hyperref}

\newtheorem{theorem}{Theorem}[section]
\newtheorem{lemma}[theorem]{Lemma}
\newtheorem{proposition}[theorem]{Proposition}
\newtheorem{corollary}[theorem]{Corollary}
\theoremstyle{definition}
\newtheorem{definition}[theorem]{Definition}
\newtheorem{remark}[theorem]{Remark}

\newcommand{\cP}{\mathcal P}
\newcommand{\cN}{\mathcal N}
\newcommand{\eps}{\varepsilon}

\title{A Boundary-Word Classification of Losing Positions in CRIM}
\author{Zijian Zeng}
\address{Institute of Computer Science and Digital Innovation, UCSI University,
Kuala Lumpur 56000, Malaysia}
\email{zijianzeng@foxmail.com}
\email{1002266693@ucsiuniversity.edu.my}
\date{August 29, 2026}
\keywords{combinatorial games, integer partitions, Young diagrams, boundary words}
\hypersetup{
  pdftitle={A Boundary-Word Classification of Losing Positions in CRIM},
  pdfauthor={Zijian Zeng}
}

\begin{document}

\begin{abstract}
Column--Row Impartial Merge (CRIM) is the normal-play impartial game in
which a move deletes an arbitrary row or column of a Young diagram and
reattaches the remaining pieces.  Ba\v{s}i\'c, Gottlieb, and Krnc conjectured
that every losing position has even absolute Dyson rank.  We prove a stronger
classification.  Encode a partition by its southeast boundary word
$0^{a_1}1^{b_1}\cdots0^{a_k}1^{b_k}$.  Let $L$ count the initial pairs with
$a_i=1$ and $b_i$ odd, let $R$ count the terminal pairs with $a_i$ odd and
$b_i=1$, and put $M=\min(k,L+R)$.  The position is losing if and only if both
the boundary length and $M$ are even.  Since the boundary length is the sum
of the number of rows and columns, the even-rank conjecture follows.  The
proof is exact and combinatorial.  Deterministic exhaustive computations,
implemented independently in three representations, are included only as
checks of the theorem and its auxiliary lemmas.
\end{abstract}

\maketitle

\section{Introduction}

Column--Row Impartial Merge (CRIM) is an impartial game on integer
partitions.  From a Young diagram, a player removes any one row or any one
column; if the removal disconnects the diagram, the remaining pieces are
reattached to form a partition.  Under normal play, a position from which the
next player loses is a $\cP$-position, and all other positions are
$\cN$-positions.

Ba\v{s}i\'c, Gottlieb, and Krnc analyzed CRIM and proved outcome
classifications for several families of partitions
\cite{BasicGottliebKrnc2026}.  Their Conjecture~1 states that every
$\cP$-position has even rank, where the rank of
$\lambda=(\lambda_1,\ldots,\lambda_r)$ is
$|r-\lambda_1|$ \cite[Sections~2.1 and~5]{BasicGottliebKrnc2026}.  This is
not the Durfee length.

We settle the conjecture by proving a complete outcome classification.  The
essential reduction replaces a partition by a binary boundary word.  In this
model, a move deletes one letter and then removes any leading 1s and trailing
0s.  Only two endpoint moves can preserve word-length parity.  Their iterated
availability is measured by an end-peeling statistic $M$.

\begin{theorem}[Boundary-word classification]\label{thm:classification}
Let $w$ be the canonical boundary word of a CRIM position, and let $M(w)$ be
the end-peeling statistic of Definition~\ref{def:M}.  Then
\[
 w\in\cP
 \quad\Longleftrightarrow\quad
 |w|\equiv0\pmod 2
 \quad\text{and}\quad
 M(w)\equiv0\pmod 2.
\]
\end{theorem}

\begin{corollary}[Even-rank conjecture]\label{cor:rank}
Every $\cP$-position in CRIM has even rank.
\end{corollary}

The converse to Corollary~\ref{cor:rank} is false: Theorem~\ref{thm:classification}
also identifies the even-rank $\cN$-positions.  The result determines only
normal-play outcomes, not all Sprague--Grundy values.

The source paper was a version-one arXiv preprint at the time of writing.
Searches of arXiv, DataCite, Crossref, OpenAlex, and the public web through
August~29, 2026 found no later proof, counterexample, or boundary-word
classification.  This is a bounded literature statement, not a claim that an
unindexed result cannot exist.

\section{Partitions and canonical boundary words}\label{sec:words}

Let $\mathbb Y$ denote the set of integer partitions, including the empty
partition $\varnothing$.  For a nonempty partition
$\lambda=(\lambda_1,\ldots,\lambda_r)$, let
\[
 v_1>v_2>\cdots>v_k>0
\]
be its distinct row lengths, and let $a_i$ be the multiplicity of $v_i$.
Set $v_{k+1}=0$ and $b_i=v_i-v_{i+1}$.  Thus $a_i,b_i$ are positive integers.

\begin{definition}\label{def:word}
The \emph{canonical boundary word} of $\lambda$ is
\[
 W(\lambda)=0^{a_1}1^{b_1}\cdots0^{a_k}1^{b_k}.
\]
We put $W(\varnothing)=\eps$.  A binary word is \emph{canonical} if it is
empty or begins with 0 and ends with 1.
\end{definition}

Every nonempty canonical word has a unique alternating-run expression as in
Definition~\ref{def:word}.  Conversely, its pairs determine the partition in
which
\[
 v_i=\sum_{j=i}^k b_j
\]
occurs $a_i$ times.  Hence $W$ is a bijection.  It also gives the parity
statistic of interest:
\begin{equation}\label{eq:length}
 |W(\lambda)|=\sum_{i=1}^k a_i+\sum_{i=1}^k b_i
 =r+\lambda_1.
\end{equation}

For any binary word $x$, let $C(x)$ remove its maximal prefix of 1s and its
maximal suffix of 0s.  If all letters are removed, set $C(x)=\eps$.

\begin{lemma}[Move correspondence]\label{lem:moves}
Under $W$, the CRIM followers of a position are precisely the distinct words
\[
 C(w\setminus x),
\]
where one occurrence $x\in\{0,1\}$ is deleted from $w$.
\end{lemma}

\begin{proof}
A row whose length is $v_i$ corresponds to one of the $a_i$ zeros in the
$i$th 0-run.  If $a_i>1$, deleting it decrements $a_i$.  If $a_i=1$ and the
run is internal, the adjacent 1-runs merge; algebraically, the gaps
$v_{i-1}-v_i$ and $v_i-v_{i+1}$ add when the sole row of length $v_i$ is
removed.  If the first 0-run disappears, its following 1-run becomes a
leading run and is stripped, recording the drop from $v_1$ to $v_2$.

The column argument is dual.  The $b_i$ columns of a fixed height correspond
to the $i$th 1-run.  Deleting a singleton internal 1-run merges its adjacent
0-runs, while deleting the last 1-run causes zero-length bottom rows to be
stripped.  These are exactly the operations described by $C$.
\end{proof}

Thus CRIM is isomorphic to the normal-play game on canonical binary words in
which a move deletes one letter and canonicalizes.  Every move strictly
shortens the word.

\section{The end-peeling statistic}

Fix a nonempty canonical word
\begin{equation}\label{eq:runs}
 w=0^{a_1}1^{b_1}\cdots0^{a_k}1^{b_k},
 \qquad a_i,b_i\ge1.
\end{equation}

\begin{definition}\label{def:M}
A run pair $(a_i,b_i)$ is \emph{A-good} if $a_i=1$ and $b_i$ is odd.  It is
\emph{B-good} if $a_i$ is odd and $b_i=1$.  Let $L(w)$ be the number of
pairs in the maximal A-good prefix, and let $R(w)$ be the number of pairs in
the maximal B-good suffix.  Define
\[
 M(w)=\min\{k,L(w)+R(w)\},
 \qquad M(\eps)=0.
\]
Equivalently, $M(w)$ is the size of the union of those two sets of pairs.
\end{definition}

The terminology comes from the parity-preserving moves.

\begin{lemma}[Even followers]\label{lem:even-followers}
If $|w|$ is even, its even-length followers are exactly the following:
\begin{enumerate}
\item if the first pair is A-good, delete its sole leading 0 and strip the
following 1-run;
\item if the last pair is B-good, delete its sole trailing 1 and strip the
preceding 0-run.
\end{enumerate}
Each such end-peel reduces $M$ by exactly one.
\end{lemma}

\begin{proof}
Unless the sole leading 0 or sole trailing 1 is deleted, canonicalization
strips nothing and exactly one letter is lost.  Deleting the sole leading 0
removes $1+b_1$ letters, preserving parity exactly when $b_1$ is odd.  The
right-end calculation removes $a_k+1$ letters and is dual.

The A-good prefix and B-good suffix are sets of intrinsically good run pairs.
An eligible end-peel removes exactly one pair from their union.  All remaining
pairs retain their A/B status.  Therefore the union size $M$ falls by one,
including when the two regions overlap.
\end{proof}

We next prove the domination statement needed for odd words.

\begin{lemma}[Odd-word domination]\label{lem:odd}
Every odd-length canonical word has a follower $u$ for which both $|u|$ and
$M(u)$ are even.
\end{lemma}

\begin{proof}
Write $L=L(w)$ and $R=R(w)$.  Every A-good or B-good pair has even total
length.  Hence $L+R<k$: otherwise their union covers every pair and $|w|$ is
even.  The nonempty interval of pairs from $L+1$ through $k-R$ will be called
the \emph{core}.  In particular,
\begin{equation}\label{eq:oddM}
 M(w)=L+R.
\end{equation}

Suppose first that the core contains at least two pairs.  If $L+R$ is odd and
$L>0$, delete a 1 from the last A-good prefix pair.  A nonsingleton 1-run
becomes even; a singleton 1-run disappears and its 0 merges with the first
core 0-run.  In either case the A-good prefix loses one pair.  A second core
pair remains as the right barrier, so the B-good suffix is unchanged.  The
new value of $M$ is $L+R-1$, which is even.  If $L=0$, then $R>0$, and the
symmetric deletion of a 0 from the first B-good suffix pair applies.

Now let $L+R$ be even, and write $(a,b)$ for the first core pair.  Except in
two cases, a one-letter deletion can leave this pair non-A-good: delete a 0,
unless $a=2$ and $b$ is odd, in which case delete a 1.  Since a second core
pair remains, the B-good suffix is unchanged, so $M$ stays even.

The first exception is $(a,b)=(2,1)$.  Delete its sole 1.  The pair merges
with the next core pair, adding 2 to that pair's 0-run.  The result is still
non-A-good, and adding 2 preserves whether the next pair is B-good; thus $M$
again stays even.

For the other exception, $a=1$ and $b$ is even.  Let $t$ be the number of
consecutive A-good pairs immediately following this core pair and set
$X=L+t+R$.  Deleting its sole 0 removes or absorbs the pair and gives
\[
 M_0=\min(k-1,X),
\]
whereas deleting one of its 1s makes the pair A-good and gives
\[
 M_1=\min(k,X+1).
\]
If $X\le k-2$, these values are $X$ and $X+1$; if $X\ge k-1$, they are
$k-1$ and $k$.  Thus they are consecutive, and one of the two legal
followers has even $M$.

It remains to consider a one-pair core $(a,b)$.  Its length is odd and it is
neither A-good nor B-good.  The following exhaustive table gives a suitable
move.  ``Reduce A'' means delete a 1 from the last A-good prefix pair;
``reduce B'' is its symmetric operation.
\begin{center}
\small
\begin{tabular}{@{}lll@{}}
\toprule
Core type & $L+R$ even & $L+R$ odd \\
\midrule
$a=1$, $b$ even
  & delete the core 0
  & delete a core 1 \\
$a$ even, $b=1$
  & delete the core 1
  & delete a core 0 \\
$a\ge3$ odd, $b$ even
  & delete a core 0
  & reduce B if $R>0$, otherwise reduce A \\
$a$ even, $b\ge3$ odd
  & delete a core 1
  & reduce A if $L>0$, otherwise reduce B \\
\bottomrule
\end{tabular}
\end{center}

For completeness, in the first two odd cases the changed core becomes good,
so all $k=L+R+1$ pairs are covered and the new $M$ is even.  In the final two
even cases, the core becomes an even--even barrier and $M$ is unchanged.  In
the final two odd cases, the indicated adjacent good region loses one pair.
If that adjacent run is a singleton, it merges into an even core run and does
not expose the opposite region.  Hence $M$ changes from $L+R$ to $L+R-1$.
All selected moves yield even word length, proving the lemma.
\end{proof}

\section{Outcome classification and rank parity}

\begin{proof}[Proof of Theorem~\ref{thm:classification}]
We use strong induction on $|w|$.  The empty word is a $\cP$-position and
satisfies the displayed criterion.  There is no nonempty canonical word of
length one.

Let $|w|$ be even.  By Lemma~\ref{lem:even-followers}, all odd-length
followers are $\cN$ by induction, and every even-length follower is an
end-peel whose $M$-value is $M(w)-1$.  If $M(w)$ is even, each such follower
has odd $M$ and is $\cN$ by induction; hence $w\in\cP$.  If $M(w)$ is odd,
an end-peel exists and leads to an even word with even $M$, which is $\cP$ by
induction.  Thus $w\in\cN$.

If $|w|$ is odd, Lemma~\ref{lem:odd} gives an even-length follower with even
$M$.  It is $\cP$ by induction, so $w\in\cN$.  This proves both directions.
\end{proof}

\begin{proof}[Proof of Corollary~\ref{cor:rank}]
By Theorem~\ref{thm:classification}, a $\cP$-position has even boundary
length.  Equation~\eqref{eq:length} and
\[
 |r-\lambda_1|\equiv r+\lambda_1\pmod2
\]
show that its rank is even.  We set the rank of the empty partition to zero;
equivalently, the nontrivial statement may be restricted to nonempty
partitions.
\end{proof}

\begin{remark}
For example, the unexplained losing partition $(4,2,2,2)$ listed in the
source paper has word $01100011$, boundary length $8$, and $M=0$, so the
classification declares it $\cP$.  The staircase $(3,2,1)$ has word
$010101$, for which $M=3$; it is therefore $\cN$ despite having even rank.
\end{remark}

\section{Independent computational verification}\label{sec:computation}

The proof above is independent of computation.  We nevertheless used three
deterministic exact implementations to attack the statement and its lemmas.

The first program enumerates partitions by size, constructs column moves by
direct column-height subtraction, and computes full Sprague--Grundy values by
mex.  The second constructs column moves by literal conjugation and computes
only the Boolean $\cP/\cN$ recursion.  On every partition of size at most 50,
the two programs produced the same losing set.

\begin{proposition}[Finite verification]
There are 1,295,971 partitions of size at most 50, including the empty
partition.  Of these, 365,327 are $\cP$-positions, and none has odd rank.
\end{proposition}

A third program works only with canonical words.  It compares direct outcome
recursion with Theorem~\ref{thm:classification}.  It separately tests
Lemma~\ref{lem:odd} and the case-selected witnesses used in its proof.  All
2,097,152 canonical words of perimeter at most 22, including the empty word,
passed with no discrepancy between the exact recursion and the predicted
outcome.

The computations used Python~3.10.12 on x86-64 Linux.  They use no randomness
and no floating-point arithmetic.  Commands, per-level counts, and
machine-readable results accompany the source distribution.

\section*{Data and code availability}

The source archive contains the three exact programs, their machine-readable
outputs, regression tests, and commands for both the fast tests and the full
finite verifications.  The programs use only the Python standard library.

\bibliographystyle{amsplain}
\bibliography{references}

\end{document}
